% =========================================================
%  INFORME DE PENTESTING - Actividad A08 - INSECUR S.A.
%  Seguridad Aplicada a Sistemas de Información
%  Compilar con: pdfLaTeX (Overleaf o TeX Live/MiKTeX)
% =========================================================
\documentclass[11pt,a4paper]{article}

\usepackage[utf8]{inputenc}
\usepackage[spanish]{babel}
\usepackage{geometry}
\geometry{left=2.2cm,right=2.2cm,top=2.5cm,bottom=2.5cm}
\usepackage[table]{xcolor}
\usepackage{tcolorbox}
\tcbuselibrary{skins,breakable}
\usepackage{enumitem}
\usepackage{tabularx}
\usepackage{booktabs}
\usepackage{longtable}
\usepackage{titlesec}
\usepackage{fancyhdr}
\usepackage{hyperref}
\usepackage{listings}
\usepackage{array}
\usepackage{amssymb}

% ----------------- Colores -----------------
\definecolor{azulmat}{RGB}{0,84,140}
\definecolor{griscaja}{RGB}{245,245,245}
\definecolor{verdesuave}{RGB}{232,244,232}
\definecolor{naranjasuave}{RGB}{255,244,230}
\definecolor{azulsuave}{RGB}{232,240,250}
\definecolor{rojosuave}{RGB}{255,235,235}

\newtcolorbox{campo}[1][Campo del Informe]{%
  enhanced, breakable, colback=white, colframe=black!60,
  fonttitle=\bfseries, title=#1,
  boxrule=0.8pt, arc=2mm, left=6pt, right=6pt, top=6pt, bottom=6pt}

\hypersetup{colorlinks=true, linkcolor=azulmat, urlcolor=azulmat, citecolor=azulmat}

\lstset{
  basicstyle=\ttfamily\small,
  backgroundcolor=\color{griscaja},
  frame=single,
  breaklines=true,
  showstringspaces=false
}

% ----------------- Encabezado -----------------
\pagestyle{fancy}
\fancyhf{}
\fancyhead[L]{\small Informe de Pentesting A08 --- INSECUR S.A.}
\fancyhead[R]{\small Seguridad Aplicada a Sistemas de Información}
\fancyfoot[C]{\thepage}
\renewcommand{\headrulewidth}{0.4pt}
\titleformat{\section}{\large\bfseries\color{azulmat}}{\thesection.}{0.6em}{}
\titleformat{\subsection}{\normalsize\bfseries\color{azulmat!80!black}}{\thesubsection.}{0.6em}{}

% ----------------- Portada -----------------
\title{%
  \rule{\textwidth}{1pt}\\[0.3cm]
  \textbf{\Large Informe de Pentesting}\\[0.2cm]
  \textbf{\Large Actividad A08 --- Hacking Ético}\\[0.2cm]
  \large Cliente: \textbf{INSECUR S.A.} --- Laboratorio: TryHackMe\\[0.2cm]
  \rule{\textwidth}{1pt}\\[0.3cm]
  \normalsize Seguridad Aplicada a Sistemas de Información\\
  Licenciatura en Sistemas de Información --- Facultad de Informática y Diseño\\
  Docente: Ing. Rodrigo Atilio Elgueta\\[0.3cm]
  \large Alumna: \textbf{Chiara Hernandez}\\
  Username de TryHackMe: \textbf{chiarahernandezp}
}
\author{}
\date{}

\begin{document}
\maketitle
\thispagestyle{fancy}
\tableofcontents
\newpage

% =========================================================
\section{Datos generales}

\begin{center}
\renewcommand{\arraystretch}{1.6}
\begin{tabularx}{\textwidth}{@{} p{5cm} X @{}}
\toprule
\textbf{Campo} & \textbf{Valor} \\
\midrule
Cliente & INSECUR S.A. \\
Contratista (alumno/a) & Chiara Hernandez \\
DNI & 45967321 \\
Username TryHackMe & chiarahernandezp \\
Email institucional & hernandezchiara@uch.edu \\
Sala de elección (Nivel 2) & [COMPLETAR] \\
Branch en GitHub & hernandez-chiara \\
Fecha de inicio & 2026-10-01 \\
Fecha de cierre & 2026-10-8 \\
\bottomrule
\end{tabularx}
\end{center}

% =========================================================
\section{Pre-pentesting}\label{sec:pre}

\begin{campo}[Alcance del trabajo]
El alcance de este trabajo se limita exclusivamente a los entornos de laboratorio
provistos por la plataforma TryHackMe (TryHackMe.com), concretamente a las máquinas
virtuales asignadas dentro de las salas \textit{Pentesting Fundamentals}, \textit{Basic
Pentesting} y la sala de elección de Nivel 2. Dentro del alcance se incluyen los
servicios expuestos por dichas máquinas (puertos TCP/UDP, servicios web, SMB, SSH,
etc.) accedidos a través de la red VPN de TryHackMe. Quedan fuera de alcance: cualquier
sistema, red o dispositivo que no pertenezca a la infraestructura de laboratorio de
TryHackMe, incluyendo redes propias, de la facultad o de terceros.
\end{campo}

\begin{campo}[Reglas de Enfrentamiento (ROE)]
\begin{itemize}[leftmargin=*]
  \item Las pruebas se realizan únicamente contra la IP asignada dinámicamente por
  TryHackMe al desplegar cada máquina de laboratorio.
  \item Está prohibido aplicar las técnicas aprendidas contra sistemas reales o
  productivos fuera del entorno de laboratorio.
  \item No se realizan ataques de denegación de servicio (DoS) ni pruebas que puedan
  degradar la disponibilidad del servicio de TryHackMe para otros usuarios.
  \item Toda la actividad se documenta en bitácora (comandos ejecutados) y evidencias
  (capturas de pantalla), conforme a los mecanismos de control de autenticidad de la
  cátedra.
  \item Las flags y contraseñas obtenidas no se publican en texto plano en el
  repositorio; se reportan mediante su hash SHA-256 (ver Anexo C).
  \item Horario de trabajo: sin restricción horaria específica, dado que el entorno es
  un laboratorio aislado sin impacto en sistemas productivos de terceros.
\end{itemize}
\end{campo}

\begin{campo}[Autorización]
Declaro que cuento con autorización expresa de la cátedra de Seguridad Aplicada a
Sistemas de Información (Docente: Ing. Rodrigo Atilio Elgueta) para realizar las
actividades de hacking ético descriptas en este informe, en el marco de la Actividad
A08, exclusivamente sobre los entornos de laboratorio de TryHackMe indicados en la
guía de la actividad.
\end{campo}

\begin{campo}[Metodología utilizada]
Se adoptó como referencia metodológica el flujo de trabajo propuesto por la guía de la
cátedra, alineado con las fases descriptas en \textbf{NIST SP 800-115} (\textit{Technical
Guide to Information Security Testing and Assessment}): Planeamiento (Pre-engagement),
Descubrimiento (Reconocimiento y Escaneo), Ataque (Explotación y Post-explotación) y
Reporte. Se eligió este marco por ser una guía técnica de referencia pública, ampliamente
utilizada en pentesting profesional, que ordena claramente las fases sin exigir
certificación previa, a diferencia de metodologías como OSSTMM. Para la fase de
reconocimiento y escaneo se utilizó además el enfoque de enumeración de servicios de
OWASP Testing Guide en lo referido a servicios web.
\end{campo}

% =========================================================
\section{Durante: Hallazgos por fase}

% Para agregar más hallazgos, copiá todo el bloque longtable y cambiá el ID (H-03, H-04...)

\begin{longtable}{@{} p{3.5cm} p{10cm} @{}}
\toprule
\textbf{Campo} & \textbf{Valor} \\
\midrule
\endhead
ID del hallazgo & H-01 \\
\midrule
Fase & $\square$ Reconocimiento $\square$ Escaneo $\boxtimes$ Explotación $\square$ Post-explotación \\
\midrule
Descripción & Dentro del home del usuario \texttt{kay} en la máquina de la sala Basic
Pentesting, el archivo de clave privada SSH \texttt{\textasciitilde/.ssh/id\_rsa} tenía
permisos \texttt{-rw-r--r--} (644), es decir, legible por cualquier usuario del sistema,
no solo por su propietario. Esto permitió, estando autenticada como el usuario
\texttt{jan}, copiar la clave privada de \texttt{kay} sin necesidad de privilegios de
root. Posteriormente, la \textit{passphrase} que protegía la clave se obtuvo mediante
ataque de diccionario offline (John the Ripper + rockyou.txt), permitiendo autenticarse
como \texttt{kay} vía SSH sin conocer su contraseña de sistema. \\
\midrule
Componente afectado & Permisos de archivo \texttt{/home/kay/.ssh/id\_rsa} \\
\midrule
Categoría (OWASP / CWE) & CWE-276 (Incorrect Default Permissions) \\
\midrule
CVE (si aplica) & N/A (error de configuración local, no vulnerabilidad de software) \\
\midrule
CVSS estimado & $\square$ Bajo $\square$ Medio $\boxtimes$ Alto $\square$ Crítico \\
\midrule
Evidencia (captura / comando) & \texttt{ls -la /home/kay/.ssh} (permisos 644 en
\texttt{id\_rsa}); \texttt{ssh2john} + \texttt{john --wordlist=rockyou.txt} (passphrase
crackeada); ver \texttt{bitacora/clase\_basic-pentesting\_04-10.txt} y
\texttt{evidencias/explotacion/}. \\
\midrule
Herramienta utilizada & OpenSSH client, ssh2john, John the Ripper \\
\midrule
Impacto en la Tríada CIA & $\boxtimes$ Confidencialidad $\boxtimes$ Integridad $\square$ Disponibilidad \\
\midrule
Salvaguarda propuesta & Restringir permisos de toda clave privada a \texttt{600}
(lectura/escritura solo para el propietario); auditar periódicamente permisos en
directorios \texttt{\textasciitilde/.ssh} de todos los usuarios; proteger claves
privadas con passphrases robustas que sigan la política de contraseñas de la
organización. \\
\midrule
Riesgo residual tras aplicar & $\boxtimes$ Bajo $\square$ Medio $\square$ Alto \\
\bottomrule
\end{longtable}

\begin{longtable}{@{} p{3.5cm} p{10cm} @{}}
\toprule
\textbf{Campo} & \textbf{Valor} \\
\midrule
\endhead
ID del hallazgo & H-02 \\
\midrule
Fase & $\square$ Reconocimiento $\square$ Escaneo $\square$ Explotación $\boxtimes$ Post-explotación \\
\midrule
Descripción & Una vez autenticada como el usuario \texttt{kay} (obtenido en H-01), se
verificaron sus privilegios con \texttt{sudo -l}. La configuración de \texttt{sudoers}
permite a \texttt{kay} ejecutar \textbf{cualquier comando como cualquier usuario,
incluyendo root, sin restricciones} (\texttt{(ALL : ALL) ALL}). Esto implica que
comprometer la cuenta de \texttt{kay} equivale a obtener control total del servidor
(por ejemplo, mediante \texttt{sudo su} o \texttt{sudo bash}). \\
\midrule
Componente afectado & Configuración \texttt{/etc/sudoers} para el usuario \texttt{kay} \\
\midrule
Categoría (OWASP / CWE) & CWE-250 (Execution with Unnecessary Privileges) \\
\midrule
CVE (si aplica) & N/A (error de configuración/política de privilegios) \\
\midrule
CVSS estimado & $\square$ Bajo $\square$ Medio $\square$ Alto $\boxtimes$ Crítico \\
\midrule
Evidencia (captura / comando) & \texttt{sudo -l} ejecutado como \texttt{kay}, salida:
\texttt{User kay may run the following commands on ip-10-66-142-44: (ALL : ALL) ALL};
ver \texttt{bitacora/clase\_basic-pentesting\_04-10.txt}. \\
\midrule
Herramienta utilizada & sudo (comando nativo del sistema) \\
\midrule
Impacto en la Tríada CIA & $\boxtimes$ Confidencialidad $\boxtimes$ Integridad $\boxtimes$ Disponibilidad \\
\midrule
Salvaguarda propuesta & Aplicar el principio de mínimo privilegio: otorgar a
\texttt{kay} permisos de sudo únicamente sobre los comandos estrictamente necesarios
para sus tareas, en lugar de acceso irrestricto (\texttt{ALL:ALL}). Revisar
periódicamente el archivo \texttt{/etc/sudoers} y los grupos con privilegios
administrativos. \\
\midrule
Riesgo residual tras aplicar & $\square$ Bajo $\boxtimes$ Medio $\square$ Alto \\
\bottomrule
\end{longtable}

% =========================================================
\section{Resumen ejecutivo de hallazgos}

\begin{center}
\renewcommand{\arraystretch}{1.4}
\begin{tabularx}{\textwidth}{@{} p{4cm} p{3cm} X @{}}
\toprule
\textbf{Severidad} & \textbf{Cantidad} & \textbf{ID de hallazgos} \\
\midrule
\cellcolor{rojosuave} Crítico & 1 & H-02 \\
\cellcolor{naranjasuave} Alto & 1 & H-01 \\
\cellcolor{verdesuave} Medio & & \\
\cellcolor{azulsuave} Bajo & & \\
\midrule
\textbf{Total} & 2 & \\
\bottomrule
\end{tabularx}
\end{center}

% =========================================================
\section{Post-pentesting: Recomendaciones}

\begin{campo}[Recomendaciones técnicas (parches y configuraciones)]
\begin{enumerate}[leftmargin=*]
  \item Corregir permisos de \texttt{/home/kay/.ssh/id\_rsa} a \texttt{600} (prioridad
  alta, asociado a H-01).
  \item Revisar y corregir la entrada de \texttt{/etc/sudoers} para \texttt{kay},
  reemplazando \texttt{ALL:ALL} por los comandos mínimos necesarios (prioridad
  crítica, asociado a H-02).
  \item Deshabilitar el directorio \texttt{/development/} del servidor web, o restringir
  su acceso, dado que exponía notas internas con información sensible (nombres de
  usuario, menciones a versiones de software).
  \item Eliminar o proteger el share SMB \texttt{Anonymous}, que permitía acceso sin
  autenticación y filtró el archivo \texttt{staff.txt} con nombres de usuarios internos.
\end{enumerate}
\end{campo}

\begin{campo}[Recomendaciones administrativas (políticas y procesos)]
\begin{enumerate}[leftmargin=*]
  \item Implementar una política de contraseñas robusta y verificar su cumplimiento
  (se detectó que la contraseña de \texttt{jan} era débil y cayó ante un diccionario
  común).
  \item Capacitar al personal sobre el uso correcto de claves SSH: generación con
  passphrase robusta y permisos restrictivos por defecto.
  \item Establecer revisiones periódicas (auditorías internas) de la configuración de
  \texttt{sudoers} y de los shares de red disponibles.
  \item Concientizar al personal sobre no dejar notas internas con credenciales,
  nombres de usuario o información sensible en recursos accesibles sin autenticación
  (web, SMB).
\end{enumerate}
\end{campo}

% =========================================================
\section{Protocolo de contingencia}

\begin{campo}[Escenario de falla]
El escenario de falla se centra en el hallazgo H-02: si un atacante externo logra
comprometer la cuenta del usuario \texttt{kay} (por ejemplo, reutilizando el vector de
H-01), obtiene de inmediato privilegios root sobre el servidor completo mediante
\texttt{sudo}. Esto podría derivar en la caída total del servicio (borrado o
modificación de archivos críticos del sistema), instalación de malware o
ransomware, o uso del servidor como pivote hacia otros sistemas de la red interna de
INSECUR S.A.
\end{campo}

\begin{campo}[Impacto sobre la continuidad del negocio]
Si el servidor comprometido aloja servicios de cara al cliente (web, archivos
compartidos vía Samba), su indisponibilidad afectaría directamente la operatoria de
INSECUR S.A.: caída del sitio web corporativo, pérdida de acceso a archivos
compartidos del equipo de desarrollo, y posible filtración de información interna
(como la expuesta en \texttt{/development/} y en el share \texttt{Anonymous}). El
tiempo de afectación dependería de si existen backups actualizados y de la rapidez de
detección del incidente; en ausencia de monitoreo activo, el compromiso podría pasar
desapercibido durante días.
\end{campo}

\begin{campo}[Plan de respuesta]
\begin{enumerate}
  \item \textbf{Activación:} Detección mediante alertas de accesos SSH no habituales,
  cambios inesperados en archivos del sistema, o reporte de indisponibilidad del
  servicio por parte de usuarios/clientes.
  \item \textbf{Contención:} Aislar el servidor de la red (desconectar o bloquear
  tráfico a nivel de firewall), revocar y rotar inmediatamente las claves SSH y
  contraseñas de todos los usuarios del sistema, en particular la de \texttt{kay}.
  \item \textbf{Recuperación:} Restaurar el servicio desde el último backup íntegro
  verificado previo al incidente; reinstalar el sistema desde una imagen limpia si se
  confirma presencia de malware o persistencia del atacante.
  \item \textbf{Comunicación:} Notificar de inmediato al responsable de sistemas/CTO de
  INSECUR S.A.; informar a clientes afectados si hubo exposición de datos; evaluar
  obligación de notificación a autoridades según la normativa de protección de datos
  vigente.
  \item \textbf{Post-mortem:} Documentar la cronología del incidente, causa raíz (en
  este caso, permisos incorrectos en \texttt{id\_rsa} y política de \texttt{sudoers}
  permisiva), y actualizar el plan de contingencia y las políticas de hardening en
  base a lo aprendido.
\end{enumerate}
\end{campo}

\begin{campo}[RTO y RPO estimados]
\textbf{RTO} (Recovery Time Objective): [COMPLETAR: estimar según criticidad real del
servicio para INSECUR S.A., p. ej. 4 horas] \\
\textbf{RPO} (Recovery Point Objective): [COMPLETAR: estimar según frecuencia de
backups disponible, p. ej. 24 horas]
\end{campo}

% =========================================================
\section{Anexos}

\begin{campo}[Anexo A --- Bitácora de comandos]
Archivos \texttt{bitacora/clase\_XX\_*.txt} incluidos en el branch.
\end{campo}

\begin{campo}[Anexo B --- Capturas de pantalla]
Carpeta \texttt{evidencias/} con subcarpetas por fase. Cada archivo nombrado como \texttt{<fase>\_<hallazgoID>.png}.
\end{campo}

\begin{campo}[Anexo C --- Hashes de las flags obtenidas]
\begin{lstlisting}
# Formato: identificador  sha256_hash
Pentesting-Fundamentals (respuesta final)
  583f30ba446f9f74d509726d25956b4b75ad965f953feffc996f6ad1382dc8d6

Basic Pentesting - contraseña de jan (hydra)
  9eb6835cc95b053818b660bd9fc9f092a0b30cef89c1f0f17f469d5caa0ae617

Basic Pentesting - passphrase clave privada de kay
  873504f0103a88d762191a358f4aa432dbbcf8af64f19dd913d7f8f925191a08

Basic Pentesting - contrasena final (pass.bak de kay)
  3ffd53ec80858f9ddcd75d69102724ab1bc1020855fff924227f618c4bfccf57

[COMPLETAR: agregar hashes de flags de la sala de elección (Nivel 2) una vez resuelta]
\end{lstlisting}
\end{campo}

\begin{campo}[Anexo D --- Declaración ética firmada]
Declaro que todas las actividades se realizaron dentro del alcance autorizado, en entornos controlados y con fines exclusivamente académicos.\\[0.4cm]
\begin{tabularx}{\textwidth}{@{} X X @{}}
\rule{6cm}{0.4pt} & \rule{6cm}{0.4pt} \\
Firma del estudiante & Fecha \\
\end{tabularx}
\end{campo}

% =========================================================
\section{Fuentes y bibliografía}

\begin{itemize}[leftmargin=*]
  \item TryHackMe. \textit{Pentesting Fundamentals}: \href{https://tryhackme.com/room/pentestingfundamentals}{sala oficial}
  \item TryHackMe. \textit{Basic Pentesting}: \href{https://tryhackme.com/room/basicpentestingjt}{sala oficial}
  \item OWASP Foundation. \textit{Web Security Testing Guide (WSTG)}: \href{https://owasp.org/www-project-web-security-testing-guide/}{sitio oficial}
  \item NIST (2008). \textit{Technical Guide to Information Security Testing and Assessment (SP 800-115)}.
\end{itemize}

\end{document}
